%% ============================================================
%%  Section: Heterogeneity by Depositor Sophistication
%%  (2026-06-12, post-JMCB revision; addresses RA4)
%%
%%  Combines the JMCB sophistication subsection (moved out of
%%  empirical_results) with three new exercises:
%%    §1 main split results (event studies + triple interactions)
%%    §2 index robustness: 2019 SOD weights
%%    §3 early/late post-period split (statistical version of the
%%       gradual-buildup visual claim)
%%    §4 subsample baselines + magnitude reconciliation
%%
%%  Cross-references consumed:
%%    fig:es_sophistication_split, tab:soph_triple_2022,
%%    tab:soph_triple_2019w, tab:soph_earlylate, tab:soph_subsample
%% ============================================================

\section{Heterogeneity by Depositor Sophistication}
\label{sec:heterogeneity}
\label{sec:results:sophistication}
\label{sec:results:heterogeneity}

The disciplinary mechanism requires depositors who are willing and able to act
on new risk information.  We test whether the response is amplified in markets
with more financially sophisticated depositor bases, splitting the sample at the
75th percentile of the deposit-weighted branch-market sophistication index of
Section~\ref{sec:data:sophistication}.  Banks at or above this threshold
are classified as operating in high-sophistication markets.

\subsection{Main Split Results}
\label{sec:heterogeneity:main}

Figure~\ref{fig:es_sophistication_split} presents event-study estimates separately
for high- and low-sophistication banks, and
Table~\ref{tab:soph_triple_2022} reports the corresponding pooled triple
interactions.  Panel~A shows that the uninsured deposit
response is concentrated in high-sophistication markets: the event-study
coefficient for that group drifts steadily downward after adoption, reaching
approximately $-1.5$ to $-2$ percentage points of assets by quarters $+9$ to $+11$,
while the low-sophistication group remains near zero throughout.  The gradual
buildup explains why the pooled triple-interaction coefficient ($\mathrm{Post}
\times \mathrm{High\,CECL} \times \mathrm{High\,Sophistication} = -0.022$,
SE $= 0.468$) is statistically insignificant: the post-period average pools early
quarters before the separation is visible with later quarters where the gap is
large.  We test this interpretation formally in
Section~\ref{sec:heterogeneity:earlylate}.  The pattern is consistent with
sophisticated depositors redirecting
balances gradually as time deposits mature, rather than breaking term contracts
early \citep{iyer2016tale}.

The sharpest statistical signal from the sophistication split appears in the price
channel.  Panel~C shows that high-CECL banks in high-sophistication markets paid
a distinctly steeper increase in interest expense after adoption; the
triple-interaction coefficient is $+0.044$ percentage points per quarter
($p = 0.020$), about 17--18 additional basis points annually on top of the
pooled effect.  Panel~D confirms that the additional funding cost translated into
larger profitability losses: the triple-interaction for ROA is $-0.076$ percentage
points per quarter ($p = 0.010$), an annualized drag of approximately 30
additional basis points.  Panel~E shows a directionally consistent but
statistically imprecise effect on NIM ($-0.031$, $p = 0.175$).  Panel~B shows
that high-sophistication banks also attracted fewer insured deposits than their
low-sophistication counterparts after adoption ($-0.965$, SE $= 0.677$,
$p = 0.154$), consistent with sophisticated depositors being more reluctant to
serve as replacement insured funding at institutions whose credit quality has
deteriorated.  Together, these results suggest a two-channel discipline
mechanism: sophisticated depositors extract higher rates immediately upon
observing the CECL disclosure, then gradually reduce balances as term instruments
mature \citep{egan2017deposit}.  We are explicit about the evidentiary
asymmetry: the price and profitability channels are statistically sharp, while
the quantity channel in high-sophistication markets is visible in the event
studies but imprecise in pooled interactions.  The remainder of this section
subjects both channels to the natural objections.

\subsection{Index Robustness: Pre-Period Branch Weights}
\label{sec:heterogeneity:index}

The baseline index weights branch ZIP-code characteristics by June 2022
deposits, raising the possibility that branch footprints or deposit
allocations had already adjusted in anticipation of the disclosure regime.
We therefore rebuild the index using June 2019 deposit weights---branch
structure fixed before both the CECL rollout and the pandemic.  The 2019- and
2022-weighted indices correlate at 0.979 across 4,261 banks, and only 2.9
percent of banks switch classification, which itself indicates that
footprints did not reshuffle in anticipation of CECL.
Table~\ref{tab:soph_triple_2019w} re-estimates the triple interactions under
the pre-period index.  If anything, the results sharpen: the insured-deposit
triple interaction becomes significant ($-1.418$, $p < 0.05$), and the
interest-expense ($+0.045$, $p < 0.05$), ROA ($-0.062$, $p < 0.05$), and NIM
($-0.039$, $p < 0.10$) interactions match the baseline.  The sophistication
heterogeneity is not an artifact of post-treatment branch weighting.

\subsection{Early Versus Late Post-Period}
\label{sec:heterogeneity:earlylate}

The event studies suggest that quantity discipline in sophisticated markets
builds over time.  Table~\ref{tab:soph_earlylate} provides the statistical
version of that visual claim, splitting the post period into early
(quarters 0--4) and late (quarter 5 onward) windows.  The price and
profitability channels strengthen with horizon exactly as the figures
suggest: the interest-expense triple interaction grows from $+0.037$
($p < 0.10$) in the early window to $+0.049$ ($p < 0.05$) in the late window,
and the ROA interaction from $-0.059$ to $-0.084$ (both $p < 0.05$).  For
uninsured quantities the late-window triple interaction is negative
($-0.21$) but remains imprecise; the insured-deposit interaction is negative
and marginally significant early ($-0.97$, $p < 0.10$) and larger but noisier
late.  The honest summary is that the data statistically establish
sophistication heterogeneity in \emph{pricing} and \emph{profitability},
while the quantity heterogeneity rests on the event-study pattern and does
not achieve conventional significance in pooled or split interactions.  We
phrase our claims accordingly.

\subsection{Subsample Baselines and Magnitude Reconciliation}
\label{sec:heterogeneity:subsample}

A triple interaction can obscure what each group's own
difference-in-differences looks like, so Table~\ref{tab:soph_subsample}
reports the baseline specification separately for high- and
low-sophistication banks.  The pattern matches the channel decomposition
above.  In high-sophistication markets, treated banks show large and
significant price and profitability effects (interest expense $+0.054$,
$p < 0.01$; ROA $-0.098$, $p < 0.01$; NIM $-0.049$, $p < 0.05$) with an
uninsured-quantity point estimate of $-0.32$ that is imprecise in the
54-treated-bank subsample.  In low-sophistication markets the quantity
estimate is $-0.42$ ($p < 0.05$, 162 treated banks) while every price and
profitability coefficient is small and insignificant.  The two subsamples
reconcile exactly with the pooled estimate: weighting the subsample
uninsured coefficients by treated-bank counts gives
$(54 \times -0.32 + 162 \times -0.42)/216 = -0.39$, against a pooled
estimate of $-0.40$ on the same sample.  The economically interesting
contrast is therefore not that sophisticated markets exhibit \emph{more}
quantity discipline---they do not, on average over the post period---but
that discipline takes a different form: in sophisticated markets banks retain
depositors by paying higher rates and absorb the cost in earnings, whereas in
unsophisticated markets the adjustment runs through quantities at largely
unchanged prices.  Both forms are discipline; they price the same disclosed
risk through different margins.
